\scanpage{11}
\mspage{9}
Cette \deleted{\unsure{la com}} opération est ``transposée'' du couple de foncteurs
\[
\begin{aligned}
\Phi_n&\xrightarrow{\alpha}\Phi_{n+m}& \deleted{\unsure{i\mapsto i}}\ i\mapsto(i-m)\\
\Phi_m&\overset{\beta}{\unsure{\text{---}|\longrightarrow}}\Phi_{n+m}& i\mapsto(i+n)
\end{aligned}
\]
où on aura $\alpha(\{n\})=\beta(\{-m\})=n-m$, et on aura
\[
\Phi_{n+m}\simeq\Phi_n\amalg_e\Phi_m,
\]
i.e. $\Phi_{n+m}$ est la ``somme amalgamée'' dans $(\mathrm{Cat})$, obtenue en ``recollant'' $\Phi_n$ et $\Phi_m$, par identification de l'``élément but'' de $\Phi_n$ avec ``l'élément source'' de $\Phi_m$.

Le foncteur composition \unsure{des|de} chemins induit
\[
\underline{\mathrm{Ch}}_{m,\Sigma}(x,y)\times\underline{\mathrm{Ch}}_{n,\Sigma}(y,z)\longrightarrow\underline{\mathrm{Ch}}_{m+n,\Sigma}(x,z)
\]
d'où en passant aux $\pi_0$
\[
\operatorname{Hom}_{\Sigma,\deleted{\unsure{\Sigma}}\,\unsure{m|n}}(x,y)\times\operatorname{Hom}_{\Sigma,\deleted{\unsure{\Sigma}}n}(y,z)\longrightarrow\operatorname{Hom}_{\Sigma,m+n}(x,z),
\]
où on \deleted{\unsure{pose}} dénote, pour $a,b\in\operatorname{Ob}X$, par
\[
\operatorname{Hom}_{\Sigma,n}(a,b)\overset{\mathrm{déf}}{=}\pi_0(\underline{\mathrm{Ch}}_{m,\Sigma}(\unsure{a},b))
\]
\deleted{l'image de $\pi_0(\underline{\unsure{\mathrm{Chem}|\mathrm{Ch}}}_{n,\Sigma}(a,b))$ dans $\pi_0(\underline{\mathrm{Chem}}_{\Sigma}(a,b))$. Notons ici que comme les foncteurs de transition
\[
\underline{\mathrm{Ch}}_{\unsure{\infty}\Sigma,n}(X;a,b)\longrightarrow\underline{\mathrm{Ch}}_{\Sigma,n}(X;\unsure{x|a},b),
\]
sont pleinement fidèles, l'}

% archive sheet number: 10
